% ======= Design ======= %
\section*{Consideration of Multiple Design Approaches}

The functional and technical requirements put forth allow for more than one design. The
approaches considered here differ principally in functional requirement F4, how spikes are
classified.

\subsection*{Approach A: Offline-Supervised Sorting (Two-Phase)}
This approach operates in two phases under different constraints. An offline phase,
run once on a workstation, sorts a reference recording via unsupervised clustering and
trains a classifier on the resulting labeled units. An online phase, running on the
SoC, performs supervised inference against this pre-trained model. The offline phase
is not subject to the power, latency, or memory constraints of the devicethe online
phase is. Its requirements therefore include the offline preparation as a
precondition: a labeled unit library and trained classifier must exist before the
device operates. [Systematic statement of the offline-phase operations and their
requirements belongs here.]

The proposed system consists of two overarching parts, or phases, under different 
constraints: one offline phase and one online phase. The offline phase runs once 
on a workstation where it sorts a reference recording and trains the classifier. The 
online phase, running on the SoC, performs inference against the resulting pre-trained 
model. Only the online phase is subject to the power, latency, and memory constraints that
motivate this thesis. The requirements described below apply exclusively to the
online phase unless otherwise noted.

\subsection*{Approach B: Online-Unsupervised Sorting (Single-Phase)}
This approach performs all sorting on the device as data arrive, using an online
unsupervised clustering algorithm. No offline phase and no reference recording are
required; the system is self-contained and needs no per-subject calibration. Because
clustering must proceed without the full data set available in advance, sorted-unit
stability and accuracy are expected to be lower than the supervised approach, which
is believed to be the central trade-off between the two designs.

\subsection*{Feature Space: Single-Channel Temporal vs.\ Multi-Channel Spatiotemporal}
Orthogonal to the supervised/unsupervised choice above is the feature space used to
discriminate units. A single channel exposes only \emph{temporal} waveform features,
whereas co-located channels (as in a tetrode or high-density probe) additionally expose
\emph{spatial} features---a spike's graded amplitude across adjacent electrodes---which
substantially improve the separation of neurons that are near-identical on any single
channel. This also sharpens the meaning of ``channels'': the bandwidth argument
(Table~2, T1) concerns \emph{independent} channels, whereas spatiotemporal sorting
requires \emph{co-located} channels that share the same neurons. Spatiotemporal
discrimination is the higher-fidelity target, consistent with dense-probe practice and
with fabricated on-chip sorters that rely on spatial features; single-channel temporal
sorting is retained as the provable minimum. The former additionally requires
probe-geometry-aware simulated data (e.g., MEArec), a dependency the single-channel case
does not impose.

% [NEXT STEP (? - potentially but have to be extremely diligent and follow the rubric and template here): 
% comparison of the two approaches against the critical technical
% requirements (accuracy, memory, calibration burden, self-containment), and
% the subsystem interactions each implies.]